\begin{corollary}[5.10.8]
\label{IV.2.s5.10.8-fr}
Sean $X$ un preesquema localmente noetheriano y $d$ un entero tal que, para
todo $x\in X$, la relación $\dim(\mathcal O_x)\geq d$ implique
$\operatorname{prof}(\mathcal O_x)\geq2$. Supongamos que $X$ es conexo;
entonces, si $X'$, $X''$ son dos componentes irreducibles distintas de $X$,
existe una sucesión $(X_i)_{0\leq i\leq n}$ de componentes irreducibles de
$X$ tal que $X_0=X'$, $X_n=X''$ y, para $1\leq i\leq n$, se tiene
$\operatorname{codim}(X_{i-1}\cap X_i,X)\leq d\setmin 1$
\textup{(se dice entonces que $X$ es \emph{conexo en codimensión} $\leq d\setmin 1$)}.
\end{corollary}

Si $Y$ es una parte cerrada de $X$ tal que
$\operatorname{codim}(Y,X)\geq d$, se tiene $\dim(\mathcal O_{X,y})\geq d$
para todo $y\in Y$ (5.1.3), y, por tanto,
$\operatorname{prof}(\mathcal O_{X,y})\geq2$ para todo $y\in Y$; de (5.10.7)
se sigue que $X\setmin Y$ es conexo. Por otra parte, para que
$\operatorname{codim}(Y,X)\geq d$, es necesario y suficiente que, para todo
$y\in Y$, exista una vecindad abierta $V$ de $y$ en $X$ tal que
$\operatorname{codim}(Y\cap V,V)\geq d$ (\textbf 0, 14.2.3). Observemos, por último,
que, si $\mathfrak F$ denota el conjunto de las partes cerradas $Y$ de $X$ de
codimensión $\geq d$, la unión de dos conjuntos de $\mathfrak F$ pertenece a
$\mathfrak F$ (\textbf 0, 14.2.5), y todo conjunto cerrado contenido en un conjunto de
$\mathfrak F$ pertenece a $\mathfrak F$; estas propiedades se expresan también
diciendo que $\mathfrak F$ es un \emph{antifiltro} de partes cerradas de $X$.
El corolario resulta entonces del siguiente lema topológico:

\begin{lemma}[5.10.8.1]
\label{IV.2.s5.10.8.1-fr}
Sean $X$ un espacio topológico conexo localmente noetheriano y $\mathfrak F$
un antifiltro de partes cerradas de $X$. Se supone que, si $Y$ es una parte
cerrada de $X$ tal que para todo $y\in Y$ existen una vecindad abierta $V$ de
$y$ en $X$ y un $Y_y\in\mathfrak F$ tales que $V\cap Y=V\cap Y_y$, entonces
$Y\in\mathfrak F$. Las condiciones siguientes son entonces equivalentes:
\begin{enumerate}
\item[\textup{a)}] Para todo $Y\in\mathfrak F$, $X\setmin Y$ es conexo.
\item[\textup{b)}] Si $X'$ y $X''$ son dos componentes irreducibles distintas de $X$,
existe una sucesión $(X_i)_{0\leq i\leq n}$ de componentes irreducibles de
$X$ tal que $X_0=X'$, $X_n=X''$ y, para $1\leq i\leq n$, se tiene
$X_{i-1}\cap X_i\notin\mathfrak F$.
\end{enumerate}
\end{lemma}

Supongamos que se verifica \emph{b)} y demostremos que $U=X\setmin Y$ es conexo para todo
$Y\in\mathfrak F$. Si $U'$ y $U''$ son dos componentes irreducibles distintas
de $U$, existen dos componentes irreducibles $X'$, $X''$ de $X$ tales que
$X'\cap U=U'$, $X''\cap U=U''$ (\textbf 0$_\mathrm I$, 2.1.6);

\oldpage[IV-2]{117}

formemos para estas dos componentes una sucesión $(X_i)$ con la propiedad
enunciada en \emph{b)} y pongamos $U_i=X_i\cap U$ para $0\leq i\leq n$; entonces
$U_i$ es una componente irreducible de $U$ (\textbf 0$_\mathrm I$, 2.1.6) y, además,
$U_i\cap U_{i-1}\neq\varnothing$ para $1\leq i\leq n$; en caso contrario, se
tendría $X_i\cap X_{i-1}\subset Y$, y por tanto
$X_i\cap X_{i-1}\in\mathfrak F$, en contradicción con la definición de los
$X_i$. Esto implica que $U$ es conexo.

Demostremos ahora que \emph{a)} implica \emph{b)}. Designemos por $Y$ la unión de la familia
$(X_\alpha\cap X_\beta)$, donde $(X_\alpha,X_\beta)$ recorre el conjunto de
los pares de componentes irreducibles \emph{distintas} de $X$ tales que
$X_\alpha\cap X_\beta\in\mathfrak F$. Para todo punto $y\in Y$ existe una
vecindad abierta $V$ de $y$ en $X$ que solo corta un número finito de
componentes irreducibles de $X$; esto muestra, por una parte, que $Y$ es
cerrado y, por otra, que $V\cap Y$ es la intersección de $V$ con un conjunto
de $\mathfrak F$; en virtud de la hipótesis impuesta a $\mathfrak F$, se tiene
$Y\in\mathfrak F$, por lo que $U=X\setmin Y$ es conexo; además, $Y$ es raro en $X$.
Sean entonces $X'$, $X''$ dos componentes irreducibles distintas de $X$, y
$U'$, $U''$ sus trazas respectivas sobre $U$; son componentes irreducibles
distintas de $U$ (\textbf 0$_\mathrm I$, 2.1.6). Ahora bien, la unión de las
componentes irreducibles $W$ de $U$ tales que existe una sucesión
$(U_i)_{0\leq i\leq n}$ de componentes irreducibles de $U$ para la cual
$U_0=U'$, $U_{i-1}\neq U_i$ y $U_{i-1}\cap U_i\neq\varnothing$ para
$1\leq i\leq n$ y $U_n=W$ es un conjunto \emph{abierto y cerrado} en $U$, puesto que
$U$ es localmente noetheriano y, por consiguiente, sus componentes
irreducibles forman una familia localmente finita de conjuntos cerrados.
Existe, por tanto, una sucesión de este tipo $(U_i)$ para la cual $U_n=U''$;
sea $X_i$ ($0\leq i\leq n$) la componente irreducible de $X$ tal que
$X_i\cap U=U_i$ (\textbf 0$_\mathrm I$, 2.1.6); como $U_{i-1}\neq U_i$, se tiene
$X_{i-1}\neq X_i$ para $1\leq i\leq n$; si se tuviera
$X_{i-1}\cap X_i\in\mathfrak F$ para un $i$ tal que $1\leq i\leq n$, se
deduciría $X_{i-1}\cap X_i\subset Y$ por definición de $Y$, de donde
$U_{i-1}\cap U_i=\varnothing$, contra la hipótesis. Esto concluye la
demostración del lema.

Se observará que la hipótesis impuesta a $X$ en (5.10.8) se verifica cuando
$X$ es un \emph{preesquema de Cohen-Macaulay y} $d\geq2$.

